\section*{Capítulo \ref{ch:b3:chromatin-epigenetics} --- Cromatina y epigenética}

\begin{solution}{exo:b3:chromatin-epigenetics:1}
Dos copias de cada una de H2A, H2B, H3 y H4 (un octámero) con
\qty{147}{bp} de ADN en $1.65$ vueltas; enlazadores de
\qtyrange{20}{60}{bp}, lo que da una repetición cercana a \qty{200}{bp}.
La nucleasa corta solo el ADN enlazador expuesto, nunca dentro de un
núcleo, así que todo fragmento es un número entero de repeticiones: una
escalera de \num{200}, \num{400} y \qty{600}{bp}, no una mancha continua.
\end{solution}

\begin{solution}{exo:b3:chromatin-epigenetics:2}
Diploide: $6.0\times 10^{9}$ bp; $6.0\times 10^{9}/190 \approx
3.2\times 10^{7}$ \omterm{def:b3:chromatin-epigenetics:nucleosome}{nucleosomas}; longitud $6.0\times 10^{9}\times
\qty{0.34}{nm} = \qty{2.0}{m}$.
\end{solution}

\begin{solution}{exo:b3:chromatin-epigenetics:3}
H3K27me3: \omterm{def:b3:chromatin-epigenetics:nucleosome}{histona} H3, lisina 27, trimetilada --- \omterm{def:b3:chromatin-epigenetics:nucleosome}{cromatina} silenciosa de
genes reprimidos durante el desarrollo, escrita por PRC2 y leída por las
proteínas con cromodominio de PRC1. H4K16ac: H4, lisina 16, acetilada ---
\omterm{def:b3:chromatin-epigenetics:nucleosome}{cromatina} activa y abierta (rompe un contacto entre \omterm{def:b3:chromatin-epigenetics:nucleosome}{nucleosomas}).
H3S10ph: H3, serina 10, fosforilada --- una marca mitótica asociada a la
condensación de los cromosomas (y, de forma transitoria, a la activación
génica por algunas señales); no es una marca estable de silenciamiento ni
de activación.
\end{solution}

\begin{solution}{exo:b3:chromatin-epigenetics:4}
El gen naranja está ligado al X y tiene un alelo naranja y otro negro; un
pelaje con manchas necesita uno de cada, en dos cromosomas X distintos,
de los cuales se silencia uno por célula. Un macho tiene un solo X y es
naranja o negro por entero. Un carey macho lleva dos cromosomas X y un Y
(XXY), y es estéril.
\end{solution}

\begin{solution}{exo:b3:chromatin-epigenetics:5}
$\qty{200}{bp}\times \qty{0.34}{nm} = \qty{68}{nm}$ por cuenta de
\qty{11}{nm}: razón $6.2$. Necesaria: $\qty{4}{cm}/\qty{4}{\micro m} =
10^{4}$. Plegamiento adicional por $10^{4}/6.2 \approx 1600$.
\end{solution}

\begin{solution}{exo:b3:chromatin-epigenetics:6}
$\mu = 0.98$, $\delta = 0.01$: $p^{*} = 0.01/0.03 = 0.33$; razón $\mu -
\delta = 0.97$; semivida $\ln 2/(-\ln 0.97) = 22.8 \approx 23$
divisiones. $\mu = 0.999$, $\delta = 0.001$: $p^{*} = 0.5$; razón
$0.998$; semivida $\ln 2 / 0.002 \approx 350$ divisiones.
\end{solution}

\begin{solution}{exo:b3:chromatin-epigenetics:7}
$A$: activo (promotor con H3K4me3 y \omterm{def:b3:chromatin-epigenetics:histone-code}{acetilación}). $B$: silencioso,
reprimido por Polycomb --- \omterm{def:b3:chromatin-epigenetics:nucleosome}{heterocromatina} facultativa, el estado de un
gen del desarrollo apagado en este linaje pero usado en otros, de modo
que $B$ es el que probablemente esté activo en otro tipo celular. $C$:
silencioso, \omterm{def:b3:chromatin-epigenetics:nucleosome}{heterocromatina} constitutiva (vía de HP1), típicamente
repeticiones o genes silenciados en todos los tipos celulares.
\end{solution}

\begin{solution}{exo:b3:chromatin-epigenetics:8}
\emph{UBE3A} solo se expresa desde el alelo materno en las neuronas. Su
deleción le vino de su padre, en el alelo que de todos modos está
silencioso, así que ella está sana; cuando la transmita, pasará a ser una
deleción materna: cada hijo tiene probabilidad $1/2$ de heredarla y
entonces tiene el síndrome de Angelman. Si la deleción le hubiera venido
de su madre, ella misma tendría el síndrome de Angelman; una portadora
sana solo puede haberla heredado por vía paterna. (Una deleción paterna
de toda la región daría un Prader--Willi; \emph{UBE3A} sola no basta.)
\end{solution}

\begin{solution}{exo:b3:chromatin-epigenetics:9}
Un X sin \emph{\omterm{def:b3:chromatin-epigenetics:x-inactivation}{Xist}} no puede inactivarse, así que el otro X se silencia
en todas las células: la inactivación es completa pero ya no aleatoria
(sesgada), y el embrión es viable. Los dos X sin \emph{\omterm{def:b3:chromatin-epigenetics:x-inactivation}{Xist}}: no hay
inactivación, dos X activos, expresión doble de los genes ligados al X, y
el embrión hembra muere pronto. Un transgén de \emph{\omterm{def:b3:chromatin-epigenetics:x-inactivation}{Xist}} en un autosoma
recubre ese autosoma en cis y silencia sus genes, lo que provoca la
pérdida de función de una copia autosómica --- el experimento que mostró
que \emph{\omterm{def:b3:chromatin-epigenetics:x-inactivation}{Xist}} basta para silenciar en cis.
\end{solution}

\begin{solution}{exo:b3:chromatin-epigenetics:10}
Las citosinas metiladas se desaminan a timina, que la reparación no
reconoce, de modo que un \omterm{def:b3:chromatin-epigenetics:methylation}{CpG} metilado muta a TpG o CpA a alta velocidad.
Solo las secuencias que están sin metilar en la \emph{línea germinal},
generación tras generación, conservan sus \omterm{def:b3:chromatin-epigenetics:methylation}{CpG}; las islas han conservado
los suyos, luego han estado sin metilar en la línea germinal durante casi
toda la historia de los vertebrados. Un gen metilado solo en el hígado
está sin metilar en el espermatozoide y el óvulo; las mutaciones que
surgen en las células hepáticas no se heredan, así que la metilación
somática no deja huella en la secuencia.
\end{solution}

\begin{solution}{exo:b3:chromatin-epigenetics:11}
En la replicación las \omterm{def:b3:chromatin-epigenetics:nucleosome}{histonas} viejas, marcadas, se reparten entre los
dos dúplex hijos y se intercalan con otras nuevas sin marcar, de modo que
cada dominio hijo está marcado a medias. HP1 unido a un \omterm{def:b3:chromatin-epigenetics:nucleosome}{nucleosoma}
marcado retenido recluta a SUV39H, que metila a los vecinos sin marcar;
mientras los \omterm{def:b3:chromatin-epigenetics:nucleosome}{nucleosomas} marcados sean lo bastante densos como para
llevar un \omterm{def:b3:chromatin-epigenetics:histone-code}{escritor} a cada hueco, el dominio se rellena antes de la
división siguiente --- la misma retroalimentación que el $\mu$ alto del
teorema. La propagación se detiene en los límites: proteínas aislantes
(CTCF), genes muy transcritos y genes de ARNt, \omterm{def:b3:chromatin-epigenetics:nucleosome}{nucleosomas} con marcas
activas incompatibles (H3K4me3, \omterm{def:b3:chromatin-epigenetics:histone-code}{acetilación}) que el \omterm{def:b3:chromatin-epigenetics:histone-code}{escritor} no puede
usar, y el punto en que los \omterm{def:b3:chromatin-epigenetics:histone-code}{borradores} (desmetilasas, recambio por HDAC)
adelantan al \omterm{def:b3:chromatin-epigenetics:histone-code}{escritor}. Prueba: recluir HP1 de forma transitoria a un gen
indicador (una fusión con un dominio de unión al ADN cuya unión se
controla con un fármaco), soltarlo y seguir el silencio del indicador y
la H3K9me3 a lo largo de las divisiones; se predice una persistencia que
depende del \omterm{def:b3:chromatin-epigenetics:histone-code}{escritor}, y su pérdida si se mutan el \omterm{def:b3:chromatin-epigenetics:histone-code}{escritor} o la
dimerización del \omterm{def:b3:chromatin-epigenetics:histone-code}{lector}, o si se coloca un aislante dentro del dominio.
\end{solution}

\begin{solution}{exo:b3:chromatin-epigenetics:12}
Hay que excluir: un ambiente compartido (la dieta y la pobreza de los
propios nietos); un efecto uterino --- la hija de la madre expuesta a la
hambruna era un feto, con sus células germinales ya presentes, de modo
que la línea germinal del nieto quedó expuesta directamente y el efecto
no es transgeneracional ---; las diferencias genéticas entre las familias
que sobrevivieron y las que no; la causalidad inversa (la propia masa
corporal altera la metilación en la sangre); las diferencias de
composición celular de la sangre; y las pruebas múltiples sobre muchos
\omterm{def:b3:chromatin-epigenetics:methylation}{CpG}. Experimento en ratones: expónganse \emph{machos} F$_0$ consanguíneos
(sin vía uterina), críense con hembras no expuestas o mediante
fecundación in vitro y transferencia de embriones, y sígase a F$_2$ y
F$_3$ por vía paterna, midiendo la metilación en el esperma de cada
generación; una marca que persiste en el esperma y el fenotipo de F$_3$,
en una cepa sin variación genética, es transmisión por la línea germinal.
\end{solution}

\begin{solution}{pb:b3:chromatin-epigenetics:1}
\textbf{1.} $6.4\times 10^{9}\times \qty{0.34}{nm} = \qty{2.2}{m}$;
$6.4\times 10^{9}/200 = 3.2\times 10^{7}$ \omterm{def:b3:chromatin-epigenetics:nucleosome}{nucleosomas}.
\textbf{2.} \omterm{def:b3:chromatin-epigenetics:nucleosome}{Histona}: $3.2\times 10^{7}\times 1.08\times 10^{5} =
3.5\times 10^{12}$ Da $= \qty{5.7}{pg}$. ADN: $6.4\times 10^{9}\times
650 = 4.2\times 10^{12}$ Da $= \qty{6.9}{pg}$. Masas comparables.
\textbf{3.} Núcleo: $\tfrac{4}{3}\pi\,(\qty{5}{\micro m})^{3} =
\qty{520}{\micro m^3}$. Partícula: $\pi\,(5.5)^{2}\times 5.5 =
\qty{520}{nm^3} = 5.2\times 10^{-7}\,\unit{\micro m^3}$; total
$3.2\times 10^{7}\times 5.2\times 10^{-7} = \qty{17}{\micro m^3}$:
alrededor del \qty{3}{\%} del núcleo.
\textbf{4.} $3.2\times 10^{7}\times \qty{11}{nm} = \qty{0.35}{m}$;
\omterm{prop:b3:chromatin-epigenetics:compaction}{razón de empaquetamiento} $2.2/0.35 \approx 6$.
\textbf{5.} $6.4\times 10^{9}/2\times 10^{5} = \num{32000}$ bucles de
\num{1000} \omterm{def:b3:chromatin-epigenetics:nucleosome}{nucleosomas} cada uno.
\textbf{6.} $0.21^{2}\times 6.4\times 10^{9} = 2.8\times 10^{8}$;
observados $5.6\times 10^{7}$, luego se ha perdido el \qty{80}{\%}, por
desaminación de la 5-metilcitosina a timina (transiciones C$\to$T sin
reparar en los \omterm{def:b3:chromatin-epigenetics:methylation}{CpG} metilados).
\textbf{7.} $p_{n+1} = 0.99\,p_{n} + 0.02\,(1 - p_{n}) = 0.97\,p_{n} +
0.02$; $p^{*} = 0.02/0.03 = 0.67$.
\textbf{8.} $p_{n} = 0.67 + 0.33\times 0.97^{n}$: $n = 10$: $0.91$;
$n = 50$: $0.74$; $n = 200$: $0.67$ (el exceso es $7\times 10^{-4}$).
\textbf{9.} $\ln 2/(-\ln 0.97) = 22.8 \approx 23$ divisiones; a una por
semana, unas \num{23} semanas --- cinco meses.
\textbf{10.} Cada replicación empareja una hebra vieja, aún metilada, con
una hebra nueva que ninguna enzima metila; las hebras viejas ni se ganan
ni se pierden, el número de hebras se duplica, y por tanto la fracción
metilada se reduce a la mitad. $0.80/2^{n} < 0.05$ exige $2^{n} > 16$:
cinco divisiones ($0.025$; con cuatro se llega exactamente al
\qty{5}{\%}).
\textbf{11.} Razón $0.9995 - 0.0005 = 0.999$; semivida $\ln
2/(-\ln 0.999) \approx 690$ divisiones, unos \num{13} años a una por
semana.
\textbf{12.} Un gen silencioso durante toda una vida de unas \num{4000}
divisiones semanales necesita una semivida de centenares de divisiones,
que solo da la retroalimentación de la pregunta~11: los \omterm{def:b3:chromatin-epigenetics:histone-code}{lectores} de la
marca (MeCP2, HP1) reclutan a sus \omterm{def:b3:chromatin-epigenetics:histone-code}{escritores} (metilasa de H3K9, DNMT3),
de modo que un dominio denso de marcas se copia con una eficiencia que
ninguna enzima aislada tiene.
\textbf{13.} Un X lleva el alelo naranja y el otro el negro; toda célula
de la piel ha silenciado un X y expresa solo el otro; las descendientes
de una célula conservan su elección, de modo que una mancha es un clon (o
un grupo de clones contiguos) que hizo la misma elección.
\textbf{14.} Las $N$ eligen el X que lleva el naranja, o todas eligen el
otro: $2\times (1/2)^{N} = 2^{1-N}$.
\textbf{15.} $2^{1-N} = 10^{-9}$: $N - 1 = \log_{2} 10^{9} = 29.9$,
$N \approx 30$ células.
\textbf{16.} $\qty{3000}{cm^2}/\qty{15}{cm^2} = 200$ manchas, muchas más
que $30$ fundadoras: durante el crecimiento las descendientes de una
fundadora se dispersan por mezcla y migración celular, de modo que un
clon forma varias islas (y las vecinas con la misma elección se funden,
lo que actúa en sentido contrario pero menos).
\textbf{17.} Un \omterm{def:b3:chromatin-epigenetics:x-inactivation}{corpúsculo de Barr} por cada X más allá del primero: XXY
uno, XXX dos, XY ninguno.
\textbf{18.} Se tolera cuando el gen tiene un homólogo funcional en el Y,
de modo que los machos también tienen dos copias (los genes
seudoautosómicos y parejas como \emph{KDM5C}/\emph{KDM5D}); importa en
las aneuploidías del X --- en el síndrome de Turner (XO) los genes que
escapan están en una sola copia, que es por lo que un individuo XO no es
simplemente una hembra normal (haploinsuficiencia de \emph{SHOX} y talla
baja), y en XXY los que escapan están sobredosificados.
\textbf{19.} $1/\num{15000} = 6.7\times 10^{-5}$; deleciones $0.70\times
6.7\times 10^{-5} = 4.7\times 10^{-5}$ (uno de cada \num{21000});
disomía materna $1.7\times 10^{-5}$ (uno de cada \num{60000}).
\textbf{20.} El Angelman es la pérdida de un solo gen, \emph{UBE3A}, así
que basta una mutación puntual en el alelo materno. El Prader--Willi es
la pérdida simultánea de varios genes de expresión paterna
(\emph{SNRPN}, el agrupamiento de ARN pequeños); una mutación puntual
inactiva uno de ellos y da como mucho un fenotipo parcial, nunca el
síndrome entero.
\textbf{21.} El \qty{20}{\%} de los hijos recibe una deleción paterna y
tiene Prader--Willi; ninguno tiene Angelman, que exige una deleción
materna.
\textbf{22.} Tres cromosomas, dos maternos y uno paterno, y se pierde uno
al azar: el paterno se pierde con probabilidad $1/3$, lo que deja una
disomía materna y un Prader--Willi; con probabilidad $2/3$ se pierde una
copia materna y el niño es normal.
\textbf{23.} Los genes paternos deberían aumentar la demanda sobre la
madre; su pérdida (Prader--Willi) debería reducirla --- lo que encaja con
la succión débil y el retraso del crecimiento del recién nacido ---,
mientras que el apetito insaciable posterior es lo contrario y es más
difícil de encajar (una lectura propuesta es que los genes paternos
empujan normalmente al niño hacia el destete). El Angelman, la pérdida de
un gen materno, debería aumentar la demanda; la succión prolongada y los
despertares nocturnos frecuentes de los niños con Angelman encajan, pero
las convulsiones y la pérdida del habla no tienen nada que ver con los
recursos.
\textbf{24.} La diana es el transcrito antisentido paterno
(\emph{UBE3A-ATS}), con un oligonucleótido antisentido que lo degrade o
bloquee su elongación: el \emph{UBE3A} paterno está intacto y solo se
silencia en cis por ese ARN, de modo que quitar el ARN restaura su
expresión. La metilación de la \omterm{def:b3:chromatin-epigenetics:imprinting}{región de control de la impronta} queda
intacta, así que \emph{SNRPN} y los demás genes improntados conservan su
patrón parental normal.
\textbf{25.} Semivida de una marca: unas $23$ divisiones sin
retroalimentación y unas $690$ con ella; el pelaje de la gata carey lo
decidió una población fundadora de unas $30$ células.
\end{solution}
